% =============================================================================
% TinyKernel — Documento Fundador
% =============================================================================
% document_id: TK-FND-00
% title: TinyKernel — Fundamentos Ontológicos de um Programa Experimental para Estruturas Causais Minimais
% subtitle: Fenômeno, realização, intervenção, evidência e minimalidade em uma ontologia operacional versionada
% version: 0.2.0
% status: BASE ONTOLÓGICA / PROGRAMA EXPERIMENTAL / PRÉ-TK-0000
% project: TinyKernel
% family: TinyLogic
% author: José Pedro Trindade
% date: 2026-09-09
% language: pt-BR
% source_primary: TK-FND-00 v0.1.0 + revisão conceitual de 2026-09-09
% source_lineage: Sister-Kernel; TinyLogicLM; TinyLogicVision
% ontology_id: TK-O
% ontology_version: 0.2.0
% epistemic_scope: ontologia de trabalho, definições operacionais, hipóteses e programa experimental
% evidence_status: pré-experimental; claims empíricos dependerão dos experimentos versionados
% implementation_authority: TK-0000 somente após preregistro próprio
% sister_kernel_authority: transferência somente por nova hipótese e revalidação
% motto: Sempre pronto. Sempre incompleto.
% =============================================================================

\documentclass[11pt,a4paper]{article}

\usepackage[T1]{fontenc}
\usepackage[utf8]{inputenc}
\usepackage[brazil]{babel}
\usepackage{lmodern}
\usepackage{microtype}
\usepackage{geometry}
\geometry{top=22mm,bottom=24mm,left=24mm,right=24mm}

\usepackage{amsmath,amssymb,mathtools}
\usepackage{booktabs,longtable,tabularx,array}
\usepackage{enumitem}
\usepackage{xcolor}
\usepackage{graphicx}
\usepackage{tikz}
\usetikzlibrary{arrows.meta,positioning,calc,fit,shapes.geometric}
\usepackage[most]{tcolorbox}
\usepackage{listings}
\usepackage{fancyhdr}
\usepackage{titlesec}
\usepackage{hyperref}

% -----------------------------------------------------------------------------
% Palette
% -----------------------------------------------------------------------------
\definecolor{TKInk}{HTML}{16202A}
\definecolor{TKBlue}{HTML}{2457A6}
\definecolor{TKTeal}{HTML}{157A7A}
\definecolor{TKGold}{HTML}{B27A14}
\definecolor{TKRed}{HTML}{A8423C}
\definecolor{TKGreen}{HTML}{3C7A57}
\definecolor{TKGray}{HTML}{F2F4F6}
\definecolor{TKMidGray}{HTML}{66727E}
\definecolor{TKDark}{HTML}{20262D}

\hypersetup{
  colorlinks=true,
  linkcolor=TKBlue,
  urlcolor=TKTeal,
  citecolor=TKBlue,
  pdftitle={TinyKernel — Fundamentos Ontológicos de um Programa Experimental para Estruturas Causais Minimais},
  pdfauthor={José Pedro Trindade},
  pdfsubject={Documento fundador e ontologia operacional do laboratório TinyKernel},
  pdfkeywords={TinyKernel, ontologia, causalidade, estruturas causais minimais, intervenção, witness, evidência, causal abstraction}
}

% -----------------------------------------------------------------------------
% Metadata commands
% -----------------------------------------------------------------------------
\newcommand{\TKProject}{TinyKernel}
\newcommand{\TKDocID}{TK-FND-00}
\newcommand{\TKVersion}{0.2.0}
\newcommand{\TKOntology}{TK-O}
\newcommand{\TKDate}{09 de setembro de 2026}
\newcommand{\TKStatus}{BASE ONTOLÓGICA / PROGRAMA EXPERIMENTAL}
\newcommand{\TKMotto}{Sempre pronto. Sempre incompleto.}

% -----------------------------------------------------------------------------
% Layout
% -----------------------------------------------------------------------------
\pagestyle{fancy}
\fancyhf{}
\lhead{\small\sffamily\TKProject}
\rhead{\small\sffamily\TKDocID\ — v\TKVersion}
\cfoot{\small\sffamily\thepage}

\titleformat{\section}
  {\Large\bfseries\sffamily\color{TKInk}}
  {\thesection}{0.75em}{}
\titleformat{\subsection}
  {\large\bfseries\sffamily\color{TKInk}}
  {\thesubsection}{0.75em}{}
\titleformat{\subsubsection}
  {\normalsize\bfseries\sffamily\color{TKInk}}
  {\thesubsubsection}{0.75em}{}

\setlength{\parindent}{0pt}
\setlength{\parskip}{0.65em}
\setlist[itemize]{topsep=0.35em,itemsep=0.25em}
\setlist[enumerate]{topsep=0.35em,itemsep=0.25em}
\setlist[description]{topsep=0.35em,itemsep=0.35em}

% -----------------------------------------------------------------------------
% Boxes
% -----------------------------------------------------------------------------
\newtcolorbox{thesisbox}[1][]{
  enhanced,
  breakable,
  colback=TKBlue!5,
  colframe=TKBlue!70,
  boxrule=0.7pt,
  arc=2mm,
  left=3mm,right=3mm,top=2.5mm,bottom=2.5mm,
  title=#1,
  fonttitle=\bfseries\sffamily
}

\newtcolorbox{ontologybox}[1][]{
  enhanced,
  breakable,
  colback=TKTeal!5,
  colframe=TKTeal!70,
  boxrule=0.7pt,
  arc=2mm,
  left=3mm,right=3mm,top=2.5mm,bottom=2.5mm,
  title=#1,
  fonttitle=\bfseries\sffamily
}

\newtcolorbox{warningbox}[1][]{
  enhanced,
  breakable,
  colback=TKGold!7,
  colframe=TKGold!70,
  boxrule=0.7pt,
  arc=2mm,
  left=3mm,right=3mm,top=2.5mm,bottom=2.5mm,
  title=#1,
  fonttitle=\bfseries\sffamily
}

\newtcolorbox{falsifybox}[1][]{
  enhanced,
  breakable,
  colback=TKRed!5,
  colframe=TKRed!70,
  boxrule=0.7pt,
  arc=2mm,
  left=3mm,right=3mm,top=2.5mm,bottom=2.5mm,
  title=#1,
  fonttitle=\bfseries\sffamily
}

\newtcolorbox{markerbox}[1][]{
  enhanced,
  breakable,
  colback=TKGray,
  colframe=TKMidGray!65,
  boxrule=0.6pt,
  arc=1.5mm,
  left=3mm,right=3mm,top=2mm,bottom=2mm,
  title=#1,
  fonttitle=\bfseries\sffamily
}

% -----------------------------------------------------------------------------
% Listings
% -----------------------------------------------------------------------------
\lstdefinestyle{tk}{
  basicstyle=\ttfamily\small,
  backgroundcolor=\color{TKGray},
  frame=single,
  rulecolor=\color{TKMidGray!40},
  framerule=0.4pt,
  xleftmargin=0.5em,
  framexleftmargin=0.5em,
  breaklines=true,
  showstringspaces=false,
  columns=fullflexible
}
\lstset{style=tk}

% =============================================================================
\begin{document}

% =============================================================================
% Title page
% =============================================================================
\begin{titlepage}
\pagecolor{TKDark}
\color{white}

\vspace*{1.5cm}

{\sffamily\small
\textbf{\TKProject}\par
Família experimental TinyLogic\par}

\vspace{1.8cm}

{\sffamily\bfseries\fontsize{27}{32}\selectfont
Fundamentos Ontológicos de um Programa Experimental para Estruturas Causais Minimais\par}

\vspace{0.6cm}

{\sffamily\fontsize{14}{19}\selectfont\color{white!78}
Fenômeno, realização, intervenção, evidência e minimalidade em uma ontologia operacional versionada\par}

\vspace{1.2cm}

\begin{tcolorbox}[
	enhanced,
	colback=white!8!TKDark,
	colframe=white!28!TKDark,
	boxrule=0.6pt,
	arc=2mm,
	left=4mm,right=4mm,top=3mm,bottom=3mm
	]
\color{white}
\textbf{Proposição de partida.}
TinyKernel investiga estruturas causais minimais por meio de uma ontologia operacional
que torna explícitos o fenômeno, suas realizações, as intervenções admissíveis, o regime
de observação, as evidências produzidas e os critérios pelos quais preservação,
necessidade, minimalidade e equivalência podem ser afirmadas.
\end{tcolorbox}

\vfill

{\sffamily
\begin{tabular}{@{}ll@{}}
Documento & \TKDocID \\
Versão & \TKVersion \\
Ontologia & \TKOntology\ v\TKVersion \\
Status & \TKStatus \\
Autor & José Pedro Trindade \\
Data & \TKDate \\
\end{tabular}
}

\vspace{1.1cm}

{\sffamily\bfseries\color{white!92}\TKMotto\par}

\end{titlepage}
\nopagecolor
\color{TKInk}

% =============================================================================
\section*{Metadados, escopo e autoridade}
\addcontentsline{toc}{section}{Metadados, escopo e autoridade}

\begin{tabularx}{\textwidth}{>{\bfseries}p{4.0cm}X}
\toprule
Projeto & TinyKernel \\
Documento & TK-FND-00 \\
Versão & 0.2.0 \\
Ontologia de trabalho & TK-O v0.2.0 \\
Estado & Base ontológica, programa experimental e pré-registro conceitual anterior a TK-0000 \\
Fonte primária & TK-FND-00 v0.1.0 e revisão conceitual consolidada em 09/09/2026 \\
Linhagem & Sister-Kernel; TinyLogicLM; TinyLogicVision \\
Objeto científico & Estruturas causais minimais relativas a fenômenos, contextos, regimes de intervenção e observação explicitados \\
Unidade de investigação & Relações entre fenômeno, realização, transformação, observação, evidência e preservação \\
Status epistemológico & Ontologia de trabalho, definições operacionais, hipóteses e programa de experimentação \\
Autoridade experimental imediata & Especificação e preregistro de TK-0000 \\
Transferência a outros projetos & Nova hipótese, novo contexto e revalidação própria \\
Princípio operacional & \textbf{READY $\neq$ COMPLETE} \\
\bottomrule
\end{tabularx}

\begin{ontologybox}[Leitura do documento]
Este documento organiza o TinyKernel como uma investigação ontologicamente explícita.
Cada termo central possui um papel definido; cada relação inferida possui um regime de
evidência; cada alegação possui uma força declarada. A ontologia é versionada junto com
o programa experimental para permitir que mudanças conceituais permaneçam rastreáveis.
\end{ontologybox}

\tableofcontents
\newpage

% =============================================================================
\section{Propósito científico}

TinyKernel é um laboratório experimental dedicado à investigação de
\textbf{estruturas causais minimais}. Seu problema central é identificar, para um
fenômeno especificado, quais organizações materiais ou computacionais preservam esse
fenômeno sob um contexto e um conjunto explícito de intervenções, e quais dessas
realizações ocupam posições minimais em uma relação de redução declarada.

A pergunta de pesquisa fundadora é:

\begin{thesisbox}[Pergunta científica]
\begin{center}
\large\bfseries
Sob quais condições uma realização pode ser sustentada experimentalmente como
causalmente minimal para a preservação de um fenômeno especificado?
\end{center}
\end{thesisbox}

A hipótese de trabalho é que parte dessa pergunta pode ser tornada empiricamente
tratável quando o experimento declara, antes da intervenção:

\begin{enumerate}
  \item o fenômeno e seu perfil constitutivo;
  \item o contexto em que a alegação vale;
  \item o espaço de realizações candidatas;
  \item as intervenções admissíveis;
  \item a relação de redução usada para comparar realizações;
  \item o regime de observação e os witnesses;
  \item a regra de adjudicação que liga evidência a claims.
\end{enumerate}

Esse arranjo transforma ``minimalidade'' de uma impressão arquitetural em uma
propriedade investigável dentro de um espaço experimental delimitado.

% =============================================================================
\section{Linhagem experimental}

\subsection{TinyLogicLM: mudança persistente como trajetória causal}

TinyLogicLM tornou observável uma cadeia completa de transformação em escala humana:

\begin{lstlisting}
simbolo
-> token
-> representacao numerica
-> transformacao
-> resultado
-> erro
-> atribuicao do erro
-> atualizacao
-> comportamento posterior
\end{lstlisting}

A contribuição para TinyKernel é a percepção de que capacidades podem depender do
\textbf{fechamento de uma trajetória causal}, e não apenas da presença de componentes.
Uma abstração candidata é:

\[
\text{estado}_t
\rightarrow
\text{ação}
\rightarrow
\text{consequência}
\rightarrow
\text{diferença}
\rightarrow
\text{atualização}
\rightarrow
\text{estado}_{t+1}.
\]

A pergunta experimental correspondente é quais relações dessa trajetória são
necessárias para que uma experiência produza mudança persistente de comportamento.

\subsection{TinyLogicVision: preservação de distinções na relação com o mundo}

TinyLogicVision acrescentou uma segunda exigência: uma inferência só pode ser
interpretada adequadamente quando a cadeia entre mundo observado, suporte,
representação e decisão permanece reconstruível.

\begin{lstlisting}
fenomeno
-> observacao
-> amostragem
-> suporte
-> representacao
-> entrada
-> transformacao
-> decisao
-> projecao
\end{lstlisting}

Dessa prática emergiram distinções como:

\begin{lstlisting}
SUPPORT != DECISION POINT
DECISION POINT != DISPLAY CELL
DISPLAY CELL != H3 CELL
RGB PREVIEW != MULTISPECTRAL TENSOR
H3 != RASTER
\end{lstlisting}

TinyKernel generaliza essa experiência por meio da noção de
\textbf{distinção constitutiva}: uma distinção integra o fenômeno quando sua fusão
altera uma dependência, capacidade discriminativa, trajetória ou atribuição exigida
pela especificação do fenômeno.

% =============================================================================
\section{Enquadramento conceitual e literatura de referência}

A base do TinyKernel se beneficia de cinco tradições próximas, usadas aqui como
instrumentos de precisão conceitual.

\subsection{Intervenção e causalidade}

Modelos causais estruturais e abordagens intervencionistas tratam intervenções e
contrafactuais como ferramentas para distinguir dependências causais de mera
associação \cite{pearl2009,woodward2003}. TinyKernel adota esse compromisso
metodológico em escala experimental: uma alegação de necessidade ganha força quando
mudanças controladas em uma realização produzem mudanças previstas no fenômeno ou em
seu perfil constitutivo.

\subsection{Mecanismos e organização}

A literatura mecanicista descreve mecanismos por entidades e atividades organizadas
de maneira produtiva \cite{machamer2000}. TinyKernel toma essa organização como
objeto de redução: componentes, relações, ordem, estado e fronteiras podem ser
intervencionados separadamente para determinar qual estrutura permanece necessária.

\subsection{Abstração e consistência causal entre realizações}

Trabalhos sobre abstração causal mostram que modelos em níveis diferentes podem ser
comparados por sua consistência frente a intervenções \cite{rubenstein2017,beckers2019}.
Em redes neurais, intervenções de intercâmbio também têm sido usadas para testar se
representações internas realizam estruturas causais interpretáveis
\cite{geiger2021}. TinyKernel incorpora essa perspectiva para comparar realizações
mínimas estruturalmente distintas.

\subsection{Minimalidade causal em descoberta causal}

Na literatura de descoberta causal, a expressão \emph{Causal Minimality} possui uso
técnico ligado à relação entre uma distribuição e a estrutura gráfica que a torna
Markov \cite{spirtes2001,peters2017}. Neste documento, a expressão
\textbf{minimalidade causal de realização} designa uma propriedade diferente e
explicitamente operacional: a minimalidade de uma realização dentro de uma relação de
redução declarada, condicionada à preservação de um fenômeno sob um protocolo.

\subsection{Ontologia formal e governança de conceitos}

Ontologias contemporâneas tratam escopo, definições, relações, axiomatização,
versionamento e interoperabilidade como elementos explícitos. A ISO/IEC 21838-2:2021,
por exemplo, formaliza uma top-level ontology em OWL~2 e Common Logic
\cite{iso21838bfo}; iniciativas como a OBO Foundry operacionalizam princípios de
escopo, definições, relações e versionamento \cite{obofoundry2021}. TinyKernel usa
esses princípios como referência de governança conceitual, mantendo TK-O como uma
ontologia de domínio própria e evolutiva.

\begin{thesisbox}[Posição do TinyKernel]
TinyKernel combina ontologia explícita, intervenção controlada, redução ordenada e
comparação entre realizações para investigar a estrutura necessária à preservação de
fenômenos. A contribuição pretendida está no método experimental de redução e na
rastreabilidade entre ontologia, intervenção, evidência e claim.
\end{thesisbox}

% =============================================================================
\section{Ontologia de trabalho TK-O v0.2.0}

A ontologia de trabalho organiza o domínio em quatro planos conectados:

\begin{center}
\begin{tikzpicture}[
  node distance=8mm,
  box/.style={draw, rounded corners, align=center, minimum width=11.0cm,
              minimum height=10mm, fill=TKGray},
  arr/.style={-{Latex[length=2.4mm]}, thick, draw=TKBlue}
]
\node[box, fill=TKTeal!8] (onto) {\textbf{Plano ontológico}\\fenômeno -- contexto -- realização -- estrutura -- trajetória};
\node[box, below=of onto, fill=TKGold!8] (exp) {\textbf{Plano experimental}\\intervenção -- transformação -- comparação -- redução};
\node[box, below=of exp, fill=TKBlue!7] (epi) {\textbf{Plano epistemológico}\\observação -- witness -- evidência -- adjudicação -- claim};
\node[box, below=of epi, fill=TKGreen!8] (inf) {\textbf{Plano inferencial}\\suficiência -- necessidade -- irredutibilidade -- minimalidade -- equivalência};
\draw[arr] (onto) -- (exp);
\draw[arr] (exp) -- (epi);
\draw[arr] (epi) -- (inf);
\end{tikzpicture}
\end{center}

\subsection{Tipos ontológicos fundamentais}

\begin{longtable}{>{\bfseries}p{3.0cm}p{4.0cm}p{7.0cm}}
\toprule
ID & Tipo & Definição de trabalho \\
\midrule
\endhead
TKO:P & Fenômeno &
Padrão distinguível de organização, capacidade, relação ou trajetória cuja
preservação constitui a pergunta experimental. \\

TKO:C & Contexto &
Conjunto explicitado de condições materiais, temporais, ambientais, computacionais e
observacionais que delimitam a validade de uma alegação sobre o fenômeno. \\

TKO:$\Phi$ & Perfil constitutivo &
Especificação das distinções, relações e restrições temporais que tornam o fenômeno
experimentalmente identificável dentro de um contexto. \\

TKO:R & Realização &
Configuração concreta de entidades, estados, relações e processos que pode instanciar
o perfil constitutivo de um fenômeno. \\

TKO:$\Sigma$ & Estrutura da realização &
Descrição das partes, relações, dependências, ordem e fronteiras relevantes de uma
realização em uma granularidade declarada. \\

TKO:I & Intervenção &
Transformação controlada aplicada a uma realização ou a uma relação selecionada para
produzir uma comparação causalmente informativa. \\

TKO:O & Observação &
Registro produzido pelo aparato ao acompanhar uma realização em uma condição
experimental. \\

TKO:W & Witness &
Procedimento explicitado que extrai, testa ou classifica propriedades relevantes de
uma ou mais observações. \\

TKO:E & Evidência &
Resultado versionado de um witness acompanhado de contexto, intervenção, observações e
proveniência suficientes para sustentar uma inferência delimitada. \\

TKO:A & Regra de adjudicação &
Regra preregistrada que relaciona evidência a uma classificação experimental ou a um
nível de claim. \\

TKO:Q & Claim &
Alegação inferencial explicitamente vinculada às evidências, ao contexto e à força de
inferência que a sustentam. \\
\bottomrule
\end{longtable}

\subsection{Perfil constitutivo do fenômeno}

O fenômeno é descrito por um perfil constitutivo:

\[
\Phi(P)=\left\langle
\mathcal{D}_P,
\mathcal{L}_P,
\mathcal{T}_P
\right\rangle,
\]

onde:

\begin{description}
  \item[$\mathcal{D}_P$ — distinções:] estados, categorias ou condições que precisam
  permanecer distinguíveis para a pergunta científica;
  \item[$\mathcal{L}_P$ — relações:] dependências, acoplamentos ou vínculos cuja
  presença integra a organização investigada;
  \item[$\mathcal{T}_P$ — restrições temporais:] ordem, persistência, precedência,
  duração ou trajetória necessárias à identidade experimental do fenômeno.
\end{description}

O contexto complementa o perfil:

\[
(P,C,\Phi(P))
\]

é a unidade mínima de especificação de uma pergunta sobre preservação.

\begin{ontologybox}[Critério ontológico de identidade experimental]
Duas observações pertencem ao mesmo fenômeno experimental quando satisfazem o mesmo
perfil constitutivo sob o contexto declarado e segundo uma regra de adjudicação
previamente registrada. Esse critério é versionado junto com o experimento.
\end{ontologybox}

% =============================================================================
\section{Relações centrais da ontologia}

TK-O trata relações como elementos de primeira classe. As relações abaixo formam a
espinha dorsal do programa experimental.

\begin{longtable}{>{\ttfamily}p{5.2cm}p{8.9cm}}
\toprule
Relação & Leitura \\
\midrule
\endhead
specified\_by$(P,\Phi)$ &
O fenômeno $P$ possui o perfil constitutivo $\Phi$. \\

situated\_in$(P,C)$ &
A alegação sobre $P$ é situada no contexto $C$. \\

candidate$(R,P,C)$ &
$R$ é uma realização candidata para $P$ em $C$. \\

has\_structure$(R,\Sigma)$ &
$\Sigma$ descreve a estrutura relevante de $R$ na granularidade adotada. \\

transforms$(I,R,R')$ &
A intervenção $I$ transforma $R$ em $R'$. \\

observes$(W,R,I,C,O)$ &
O witness $W$ participa da obtenção da observação $O$ em uma condição experimental. \\

produces$(W,O,E)$ &
A aplicação de $W$ a $O$ produz a evidência $E$. \\

supports$(E,Q)$ &
A evidência $E$ sustenta o claim $Q$ na força inferencial declarada. \\

reduces$(R',R;\mathcal I,\Gamma)$ &
$R'$ é uma redução admissível de $R$ segundo a família de intervenções $\mathcal I$
e o critério de ordem $\Gamma$. \\

preserves$(R,P,C,W,A)$ &
$R$ preserva $P$ sob o contexto, witnesses e regra de adjudicação declarados. \\

causal\_equivalent$(R_1,R_2)$ &
$R_1$ e $R_2$ apresentam a mesma assinatura intervencional relevante para $\Phi$ sob
a família de provas $\mathcal J$. \\
\bottomrule
\end{longtable}

Essas relações permitem que um claim científico seja reconstruído como um caminho
explícito de proveniência:

\begin{lstlisting}
phenomenon
-> constitutive_profile
-> context
-> candidate_realization
-> intervention
-> observation
-> witness
-> evidence
-> adjudication
-> claim
\end{lstlisting}

% =============================================================================
\section{Observação, evidência e adjudicação}

Para uma realização $R$, um contexto $C$ e uma intervenção $I$, o aparato produz uma
observação:

\[
O = \operatorname{Obs}(R,C,I).
\]

Um conjunto de witnesses $W$ transforma observações em evidência estruturada:

\[
E = W(O;\Phi(P),C,I).
\]

Uma regra de adjudicação $A$ associa a evidência a uma classificação experimental:

\[
A(E,\Phi(P)) \in
\{\texttt{PRESERVED},\texttt{BROKEN},\texttt{INCONCLUSIVE},
\texttt{WITNESS\_COMPROMISED}\}.
\]

A função de preservação é então definida operacionalmente por:

\[
\operatorname{Pres}(R;P,C,W,A)=1
\quad\Longleftrightarrow\quad
A(E,\Phi(P))=\texttt{PRESERVED}.
\]

Essa definição vincula preservação a um protocolo público e versionado. O fenômeno,
a observação e a decisão experimental ocupam posições ontológicas distintas, mas
permanecem ligados por relações rastreáveis.

\subsection{Dimensões de preservação}

O perfil de witnesses pode observar diferentes dimensões:

\begin{longtable}{>{\bfseries}p{3.2cm}p{10.8cm}}
\toprule
Dimensão & Pergunta experimental \\
\midrule
\endhead
Operacional &
A realização continua produzindo o comportamento ou artefato requerido? \\

Causal &
A saída continua dependente das relações causais especificadas no protocolo? \\

Discriminativa &
As distinções constitutivas continuam preservadas e capazes de separar estados que a
pergunta científica exige distinguir? \\

Observacional &
O aparato mantém capacidade de detectar e medir as propriedades usadas na adjudicação? \\

Temporal &
Persistência, ordem e trajetória continuam satisfazendo as restrições temporais do
perfil constitutivo? \\
\bottomrule
\end{longtable}

A ruptura experimental é representada como a falha de uma ou mais dessas dimensões,
com classificação específica e proveniência preservada.

% =============================================================================
\section{Intervenção e relação de redução}

A operação fundamental de TinyKernel é uma transformação controlada:

\[
I:R\mapsto R'.
\]

Uma intervenção pode remover, substituir, fundir, reordenar, desativar ou alterar
uma relação, estado, fronteira ou etapa temporal. O protocolo declara uma família de
intervenções admissíveis $\mathcal I$.

Para falar em minimalidade, o experimento também declara um critério de comparação
$\Gamma$. Definimos:

\[
R' \prec_{\mathcal I,\Gamma} R
\]

quando existe uma intervenção admissível que transforma $R$ em $R'$ e $R'$ ocupa
posição estritamente inferior a $R$ segundo $\Gamma$.

$\Gamma$ pode expressar, conforme o experimento:

\begin{itemize}
  \item inclusão estrutural;
  \item número de relações causais ativas;
  \item número de estados ou processos independentes;
  \item profundidade ou extensão de uma trajetória;
  \item conjunto de distinções ontológicas necessárias;
  \item outra ordem parcial definida e justificada antes da busca.
\end{itemize}

\begin{thesisbox}[Minimalidade como propriedade ordenada]
Minimalidade é definida em relação a uma ordem explícita sobre realizações. Isso
permite comparar estruturas sem identificar ``menor'' com linhas de código, classes,
bytes ou preferência arquitetural.
\end{thesisbox}

% =============================================================================
\section{Suficiência, necessidade, irredutibilidade e minimalidade}

Defina o conjunto de realizações preservadoras:

\[
\mathcal S_{P,C,W,A}
=
\left\{
R\in\mathcal R:
\operatorname{Pres}(R;P,C,W,A)=1
\right\}.
\]

Uma realização $R$ é \textbf{suficiente sob o protocolo} quando:

\[
R\in\mathcal S_{P,C,W,A}.
\]

Uma relação, parte ou distinção $e$ de $R$ possui \textbf{necessidade observada}
quando uma intervenção comparável sobre $e$ produz uma realização $R'$ que perde a
preservação:

\[
R' = I_e(R),
\qquad
\operatorname{Pres}(R)=1,
\qquad
\operatorname{Pres}(R')=0.
\]

Uma realização é \textbf{irredutível relativa à família de intervenções} quando toda
redução imediatamente admissível rompe a preservação:

\[
\forall R'
\left(
R'\prec_{\mathcal I,\Gamma}R
\land
\operatorname{Cover}(R',R)
\right)
\Rightarrow
R'\notin\mathcal S.
\]

O conjunto de realizações \textbf{minimais} é:

\[
\boxed{
\mathcal M(P,C,W,A,\mathcal I,\Gamma)
=
\left\{
R\in\mathcal S:
\nexists R'\in\mathcal S
\text{ com }
R'\prec_{\mathcal I,\Gamma}R
\right\}.
}
\]

Essa formulação é central porque $\mathcal M$ pode conter uma, várias ou nenhuma
realização sustentada pelo espaço de busca examinado.

\subsection{Necessidade conjunta e redundância}

Para um subconjunto $S$ de elementos ou relações:

\[
I^{-}_{S}(R)=R\setminus S.
\]

Se intervenções individuais preservam o fenômeno e uma intervenção conjunta o rompe,
a necessidade é registrada como propriedade de uma combinação:

\[
\operatorname{Nec}(S\mid R;P,C,W,A)=1.
\]

A ontologia registra, portanto, necessidade elementar, necessidade conjunta,
redundância e substituibilidade como relações experimentalmente distintas.

% =============================================================================
\section{Assinatura intervencional e equivalência causal}

Realizações estruturalmente diferentes podem produzir o mesmo fenômeno. A comparação
entre elas deve envolver o comportamento sob intervenções relevantes, e não apenas a
saída basal.

Para uma família de provas $\mathcal J$, definimos uma assinatura intervencional:

\[
\sigma_{\mathcal J}(R;\Phi,C)
=
\left\langle
A\bigl(E_j(R),\Phi\bigr)
\right\rangle_{j\in\mathcal J}.
\]

Duas realizações são experimentalmente equivalentes no escopo quando suas assinaturas
coincidem sob uma correspondência declarada de intervenções e observáveis:

\[
R_1\sim_{\Phi,C,\mathcal J}R_2.
\]

A equivalência é sempre indexada ao fenômeno, ao contexto e à família de provas.
Ela permite comparar implementações diferentes sem perder o vínculo com o que foi
efetivamente testado.

\subsection{Hipótese de Kernel como classe de equivalência}

A partir do conjunto de realizações minimais, TinyKernel introduz a seguinte hipótese
de trabalho:

\begin{thesisbox}[Hipótese K]
Um \textbf{Kernel candidato} pode ser representado por uma classe de equivalência de
realizações minimais que compartilham uma assinatura intervencional relevante para o
mesmo perfil constitutivo.
\end{thesisbox}

Formalmente, o espaço de Kernels candidatos pode ser escrito como:

\[
\boxed{
\mathcal K^{\ast}(P,C)
=
\mathcal M(P,C,W,A,\mathcal I,\Gamma)
\big/
\sim_{\Phi,C,\mathcal J}
}
\]

Essa construção desloca a atenção da implementação concreta para invariâncias
intervencionais que sobrevivem entre realizações minimais.

A unidade científica de interesse passa a ser a sequência:

\[
\boxed{
P
\rightarrow
\Phi
\rightarrow
\mathcal S
\rightarrow
\mathcal M
\rightarrow
\mathcal K^{\ast}
}
\]

Cada seta corresponde a uma operação explicitável, testável e versionável.

% =============================================================================
\section{Ciclo experimental TinyKernel}

O ciclo básico conecta ontologia e experimento:

\begin{center}
\begin{tikzpicture}[
  node distance=8mm,
  box/.style={draw, rounded corners, align=center, minimum width=5.2cm,
              minimum height=9mm, fill=TKGray},
  arr/.style={-{Latex[length=2.4mm]}, thick, draw=TKBlue}
]
\node[box, fill=TKTeal!9] (phi) {\textbf{Especificar}\\$P,C,\Phi$};
\node[box, below=of phi] (r) {\textbf{Materializar}\\realização candidata $R$};
\node[box, below=of r] (base) {\textbf{Caracterizar baseline}\\observação + witnesses};
\node[box, below=of base, fill=TKGold!9] (i) {\textbf{Intervir}\\$I:R\mapsto R'$};
\node[box, below=of i] (e) {\textbf{Produzir evidência}\\$O\rightarrow W\rightarrow E$};
\node[box, below=of e] (a) {\textbf{Adjudicar}\\$E\rightarrow A\rightarrow Q$};
\node[box, below=of a, fill=TKGreen!9] (m) {\textbf{Atualizar espaço de hipóteses}\\$\mathcal S,\mathcal M,\mathcal K^{\ast}$};

\draw[arr] (phi) -- (r);
\draw[arr] (r) -- (base);
\draw[arr] (base) -- (i);
\draw[arr] (i) -- (e);
\draw[arr] (e) -- (a);
\draw[arr] (a) -- (m);
\draw[arr] (m.east) .. controls +(2.0,0) and +(2.0,0) .. (phi.east);
\end{tikzpicture}
\end{center}

O ciclo admite materialização e redução como operações complementares:

\[
\boxed{
\text{materializar}
\;\longleftrightarrow\;
\text{intervir}
\;\longleftrightarrow\;
\text{reduzir}
}
\]

O objetivo é tornar cada passagem justificável por uma pergunta experimental, e cada
claim reconstruível a partir das evidências que o sustentam.

% =============================================================================
\section{Escada de força dos claims}

Para impedir que uma observação local seja promovida diretamente a conclusão geral,
TinyKernel classifica a força das alegações.

\begin{longtable}{>{\bfseries}p{1.4cm}p{4.1cm}p{9.0cm}}
\toprule
Nível & Claim & Evidência requerida \\
\midrule
\endhead
L0 & Especificação ontológica &
$P$, $C$, $\Phi$, tipos e relações definidos e versionados. \\

L1 & Observação &
Resultado empírico reproduzível em uma condição declarada. \\

L2 & Suficiência relativa &
Realização preserva o fenômeno segundo witnesses e adjudicação preregistrados. \\

L3 & Necessidade relativa &
Intervenção controlada em uma relação ou conjunto produz perda de preservação. \\

L4 & Irredutibilidade &
Todas as reduções imediatas relevantes examinadas rompem preservação. \\

L5 & Minimalidade relativa &
Nenhuma realização preservadora inferior foi encontrada no espaço de busca e ordem
de redução declarados. \\

L6 & Equivalência causal experimental &
Realizações minimais compartilham assinatura intervencional sob uma família de provas. \\

L7 & Robustez contextual &
A estrutura preserva a alegação em contextos adicionais previamente especificados. \\

L8 & Kernel candidato &
Classe de equivalência ou invariante causal sustentado por evidência acumulada em
L5--L7. \\
\bottomrule
\end{longtable}

Essa escada permite que o laboratório publique resultados fortes ou modestos sem
confundir os respectivos alcances.

% =============================================================================
\section{O primeiro fenômeno: persistência adaptativa}

TK-0001 usará um fenômeno deliberadamente pequeno e observável.

\begin{thesisbox}[Fenômeno TK-0001]
\textbf{Persistência adaptativa:}
um sistema sofre uma experiência $E$, altera algum estado interno em função dessa
experiência e, em uma interação posterior na qual $E$ não é reapresentada, produz um
comportamento diferente do baseline em consequência daquela alteração persistida.
\end{thesisbox}

Um perfil constitutivo inicial pode ser:

\[
\Phi_{PA}
=
\left\langle
\mathcal D_{PA},
\mathcal L_{PA},
\mathcal T_{PA}
\right\rangle
\]

com:

\begin{description}
  \item[$\mathcal D_{PA}$:] estado antes/depois; experiência presente/ausente;
  comportamento baseline/comportamento posterior;
  \item[$\mathcal L_{PA}$:] experiência influencia alteração persistida; alteração
  persistida influencia comportamento posterior;
  \item[$\mathcal T_{PA}$:] a alteração ocorre após a experiência e permanece até
  uma interação posterior sem reapresentação de $E$.
\end{description}

Uma realização inicial pode conter:

\begin{lstlisting}
estado
-> entrada
-> acao
-> feedback
-> diferenca
-> atualizacao
-> estado persistido
-> interacao posterior
\end{lstlisting}

Intervenções iniciais candidatas:

\begin{longtable}{p{4.0cm}p{4.8cm}p{5.3cm}}
\toprule
Intervenção & Previsão preregistrada & Relação investigada \\
\midrule
\endhead
Remover atualização &
A experiência deixa de alterar o estado persistido &
Dependência entre experiência e mudança persistente \\

Remover persistência &
Mudança transitória deixa de alcançar a interação posterior &
Restrição temporal de continuidade \\

Remover feedback &
A dinâmica pode mudar de classe ou permanecer adaptativa por outra relação &
Papel causal do feedback nesta realização \\

Substituir feedback por ruído &
Estado pode variar sem permanecer orientado pela consequência &
Distinção entre adaptação orientada e mutação \\

Fundir estado e ação &
Saída imediata pode persistir enquanto a trajetória explicativa se altera &
Necessidade de estado separável para persistência causal \\
\bottomrule
\end{longtable}

% =============================================================================
\section{Programa experimental inicial}

\begin{longtable}{p{2.1cm}p{4.2cm}p{7.8cm}}
\toprule
Experimento & Objeto & Pergunta \\
\midrule
\endhead
TK-0000 &
Aparato experimental &
O ciclo especificação--baseline--intervenção--observação--evidência--adjudicação é
reproduzível, determinístico e rastreável? \\

TK-0001 &
Persistência adaptativa &
Quais relações preservam o perfil constitutivo de uma mudança persistente de
comportamento após uma experiência? \\

TK-0002 &
Observabilidade &
Como intervenções no aparato de observação alteram a confiança na adjudicação sem
alterar necessariamente a realização? \\

TK-0003 &
Realizações alternativas &
Realizações estruturalmente distintas podem ocupar $\mathcal M$ e compartilhar uma
assinatura intervencional? \\

TK-0004 &
Redundância e necessidade conjunta &
Quais relações necessárias emergem apenas em intervenções combinadas? \\

TK-0005 &
Granularidade ontológica &
Como diferentes decomposições de uma mesma realização alteram a relação de redução e
o conjunto de minimais? \\

TK-0010 &
Colapso de distinções &
Quais distinções do perfil constitutivo sustentam capacidade discriminativa,
proveniência ou atribuição causal? \\

TK-0100 &
Composição federada &
Quais relações entre Participant, Role, Authority, Evidence e Decision permanecem
necessárias em fenômenos controlados de composição? \\
\bottomrule
\end{longtable}

% =============================================================================
\section{Colapso de distinções como intervenção ontológica}

Uma contribuição metodológica específica de TinyKernel é tratar a fusão de
categorias como intervenção experimental.

Considere uma distinção:

\[
d=(x_1,x_2)\in\mathcal D_P.
\]

Uma intervenção de colapso aplica um operador $\rho$ tal que:

\[
\rho(x_1)=\rho(x_2).
\]

A pergunta experimental é se essa fusão altera a preservação de $\Phi(P)$.
Quando dois estados que o perfil exige distinguir tornam-se indistinguíveis e uma
inferência relevante deixa de ser sustentada, registra-se perda de
\textbf{capacidade discriminativa constitutiva}.

Exemplos provenientes da linhagem TinyLogicVision incluem:

\begin{lstlisting}
support := decision_point
preview_rgb := multispectral_tensor
h3_cell := raster_support
display_cell := observation_support
\end{lstlisting}

A ontologia descreve essas intervenções de forma positiva: cada distinção possui uma
função no perfil constitutivo e pode ser submetida a um teste explícito de colapso.

% =============================================================================
\section{Versionamento ontológico e materialização computacional}

TK-O evolui junto com os experimentos. Cada versão deve registrar:

\begin{enumerate}
  \item identificadores estáveis para tipos e relações;
  \item definição textual operacional;
  \item domínio e alcance de cada relação;
  \item mudanças de significado entre versões;
  \item experimentos que utilizam cada termo;
  \item claims afetados por uma mudança ontológica;
  \item marcadores de investigação associados.
\end{enumerate}

Uma representação computacional futura pode mapear TK-O para OWL~2, Common Logic,
RDF/SHACL ou outra linguagem adequada. O requisito central é que a serialização
preserve os conceitos e relações versionados no registro científico.

\subsection{Forma mínima sugerida para um termo}

\begin{lstlisting}
id: TKO:Phenomenon
label: Fenomeno
definition: Padrao distinguivel de organizacao, capacidade,
            relacao ou trajetoria cuja preservacao constitui
            a pergunta experimental.
version: 0.2.0
relations:
  - specified_by -> TKO:ConstitutiveProfile
  - situated_in -> TKO:Context
used_by:
  - TK-0001
status: ACTIVE
\end{lstlisting}

Essa forma aproxima documento, ontologia e execução sem transformar a ontologia em
um artefato meramente documental.

% =============================================================================
\section{Constituição metodológica v0.2.0}

\begin{enumerate}[label=\textbf{P\arabic*.}]
  \item \textbf{Primazia do fenômeno.}
  Toda investigação começa por um fenômeno situado e por um perfil constitutivo
  versionado.

  \item \textbf{Ontologia explícita.}
  Tipos, relações, granularidade e critérios de identidade relevantes são declarados
  antes da inferência sobre minimalidade.

  \item \textbf{Intervenção como evidência de dependência.}
  Necessidade e suficiência são fortalecidas por transformações controladas e
  comparáveis.

  \item \textbf{Separação epistemológica.}
  Fenômeno, observação, witness, evidência, adjudicação e claim possuem papéis próprios
  e proveniência conectada.

  \item \textbf{Preservação multidimensional.}
  Operação, causalidade, discriminação, observabilidade e temporalidade podem compor o
  perfil de preservação.

  \item \textbf{Minimalidade ordenada.}
  Toda alegação de minimalidade declara a relação $\prec_{\mathcal I,\Gamma}$ que a
  torna significativa.

  \item \textbf{Pluralidade de realizações.}
  O espaço experimental admite múltiplos elementos minimais e múltiplas classes de
  equivalência.

  \item \textbf{Interações explícitas.}
  Dependência conjunta, redundância e substituibilidade são investigadas por
  intervenções combinadas.

  \item \textbf{Contextualidade declarada.}
  Claims carregam o contexto, o protocolo e a força inferencial que delimitam seu
  alcance.

  \item \textbf{Preregistro e história experimental.}
  Previsões, intervenções, critérios e resultados são registrados de forma imutável
  por experimento.

  \item \textbf{Evolução ontológica versionada.}
  Mudanças de conceito ou relação geram nova versão e preservam impacto sobre
  experimentos anteriores.

  \item \textbf{Transferência por revalidação.}
  Aplicação a outro domínio começa por mapear fenômeno, contexto e relações e por
  produzir evidência própria no destino.
\end{enumerate}

% =============================================================================
\section{Transferência para Sister e outros domínios}

A transferência de resultados de TinyKernel é tratada como uma nova hipótese de
correspondência ontológica.

Uma transferência deve declarar:

\begin{enumerate}
  \item o perfil constitutivo no domínio de origem;
  \item o contexto e a família de intervenções de origem;
  \item a estrutura ou classe minimal sustentada;
  \item o mapeamento ontológico proposto para o domínio de destino;
  \item as relações preservadas pelo mapeamento;
  \item os witnesses e critérios de adjudicação do destino;
  \item o experimento capaz de sustentar ou refinar a transferência.
\end{enumerate}

Para Sister, exemplos futuros incluem relações entre:

\begin{lstlisting}
Participant
Role
Authority
Evidence
Decision
Binding
Artifact
Provenance
\end{lstlisting}

O objetivo será testar quais distinções e dependências são constitutivas de fenômenos
controlados de identidade, autoridade, decisão e composição federada.

% =============================================================================
\section{Marcadores de investigação}

\begin{markerbox}[TK-MARKER-001 — Critério de identidade do fenômeno | OPEN]
Comparar perfis constitutivos alternativos para o mesmo domínio e determinar quais
mudanças configuram refinamento de $\Phi$ e quais configuram um novo fenômeno
experimental.
\end{markerbox}

\begin{markerbox}[TK-MARKER-002 — Ordem de redução | OPEN]
Investigar quais ordens parciais $\Gamma$ permitem comparar realizações sem tornar a
minimalidade um reflexo arbitrário da decomposição escolhida.
\end{markerbox}

\begin{markerbox}[TK-MARKER-003 — Granularidade ontológica | OPEN]
Determinar como funções, estados, relações, fronteiras, informação e distinções
semânticas podem ocupar unidades comparáveis de intervenção.
\end{markerbox}

\begin{markerbox}[TK-MARKER-004 — Assinatura intervencional | OPEN]
Definir condições suficientes para que duas realizações minimais sejam tratadas como
causalmente equivalentes sob uma família de provas.
\end{markerbox}

\begin{markerbox}[TK-MARKER-005 — Fenômenos temporais | OPEN]
Formalizar perfis em que identidade depende de trajetória e ordem legítima das
transições, permitindo investigar estruturas minimais de continuidade dinâmica.
\end{markerbox}

\begin{markerbox}[TK-MARKER-006 — Capacidade discriminativa | OPEN]
Relacionar colapso de distinções a perda mensurável de inferência, proveniência,
atribuição ou decisão sem reduzir semântica a mera equivalência de valores.
\end{markerbox}

\begin{markerbox}[TK-MARKER-007 — Cobertura do espaço de busca | OPEN]
Definir como registrar a extensão efetivamente explorada de $\mathcal R$ e distinguir
minimalidade sustentada no espaço examinado de alegações de cobertura mais ampla.
\end{markerbox}

\begin{markerbox}[TK-MARKER-008 — Robustez entre contextos | OPEN]
Investigar como classes de equivalência e invariantes causais se transformam quando o
contexto muda de $C_1$ para $C_2$.
\end{markerbox}

% =============================================================================
\section{Estado operacional do programa}

O primeiro estado executável deve materializar apenas o aparato necessário a TK-0000.
A estrutura inicial pode permanecer austera:

\begin{lstlisting}
TinyKernel/
|-- CMakeLists.txt
|-- README.md
|-- bin/
|   `-- tinykernel
|-- ontology/
|   `-- TK-O-0.2.0.yaml
|-- src/
|-- tests/
`-- experiments/
    |-- TK-0000/
    |   |-- question.md
    |   |-- preregistration.md
    |   `-- result.md
    `-- TK-0001/
        |-- question.md
        |-- preregistration.md
        `-- result.md
\end{lstlisting}

O primeiro executável precisa suportar somente a superfície necessária para:

\begin{lstlisting}
./bin/tinykernel verify
./bin/tinykernel experiment TK-0000
\end{lstlisting}

TK-0000 deverá produzir um caminho completo e reproduzível entre especificação,
realização, intervenção, observação, evidência e adjudicação antes de qualquer claim
de minimalidade.

% =============================================================================
\section{\texorpdfstring{READY $\neq$ COMPLETE}{READY != COMPLETE}}

A disciplina de desenvolvimento acompanha a disciplina epistemológica.

\subsection{READY}

Um estado publicável possui:

\begin{lstlisting}
HEAD compila
HEAD executa
HEAD verifica
ontologia usada pelo experimento esta versionada
experimentos publicados sao reproduziveis
witnesses e regras de adjudicacao sao explicitos
claims apontam para evidencia e escopo
\end{lstlisting}

\subsection{INCOMPLETE}

Um estado científico permanece aberto a:

\begin{lstlisting}
novos contextos
novas realizacoes
novas intervencoes
novas granularidades
novos witnesses
novas classes minimais
revisoes ontologicas versionadas
\end{lstlisting}

\begin{thesisbox}[Invariante de projeto]
Completude operacional significa que o estado atual pode ser executado e auditado.
Abertura epistemológica significa que novas evidências podem ampliar ou reorganizar a
ontologia e o espaço de hipóteses sem apagar a história experimental.
\end{thesisbox}

% =============================================================================
\section{Síntese}

TinyKernel organiza a investigação da minimalidade causal de realização como uma
sequência explícita:

\[
\boxed{
\begin{aligned}
\text{fenômeno} &\rightarrow \text{perfil constitutivo} \rightarrow \text{realizações} \rightarrow \text{intervenções} \\n&\rightarrow \text{evidências} \rightarrow \text{ordem de redução} \rightarrow \text{minimais} \\n&\rightarrow \text{equivalência causal}
\end{aligned}
}
\]

O conceito de Kernel adquire, nesse quadro, uma forma precisa e experimentalmente
progressiva: uma classe candidata de realizações minimais cuja equivalência é
sustentada por assinaturas intervencionais relevantes para o mesmo fenômeno.

A ontologia torna explícito \emph{o que existe no universo experimental}; o protocolo
determina \emph{como essas entidades e relações podem ser perturbadas}; o aparato de
observação determina \emph{que evidência é produzida}; e a escada de claims determina
\emph{quão forte pode ser a inferência resultante}.

Essa integração entre ontologia, intervenção e evidência constitui a base científica
do programa TinyKernel.

\vspace{0.8cm}

\begin{center}
{\sffamily\Large\bfseries\color{TKBlue}
Sempre pronto. Sempre incompleto.}
\end{center}

% =============================================================================
\appendix

\section{Registro de proveniência da revisão v0.2.0}

\begin{tabularx}{\textwidth}{>{\bfseries}p{4.1cm}X}
\toprule
Versão de origem & TK-FND-00 v0.1.0 \\
Data da revisão & 09/09/2026 \\
Mudança estrutural & Reorganização do documento em torno de uma ontologia operacional e
de uma cadeia formal entre fenômeno, realização, intervenção, observação, evidência e claim \\
Mudança formal & Introdução de $\Phi(P)$, da relação de redução
$\prec_{\mathcal I,\Gamma}$, do conjunto $\mathcal M$ de realizações minimais e do
quociente experimental $\mathcal K^{\ast}=\mathcal M/\sim$ \\
Mudança epistemológica & Introdução de escada L0--L8 para graduar a força dos claims \\
Mudança semântica & Reformulação de ruptura semântica como perda de capacidade
discriminativa constitutiva sob colapso de distinções \\
Mudança metodológica & Ontologia TK-O passa a ser versionada junto com experimentos e claims \\
Literatura adicionada & causalidade interventional, mecanismos, causal abstraction,
causal discovery e governança contemporânea de ontologias \\
Próximo estado & Preregistro de TK-0000 e materialização mínima de TK-O em formato legível por máquina \\
\bottomrule
\end{tabularx}

\section{Histórico de versões}

\begin{longtable}{p{2.0cm}p{2.7cm}p{9.7cm}}
\toprule
Versão & Data & Mudança \\
\midrule
\endhead
0.1.0 & 09/09/2026 &
Primeiro registro fundador: minimalidade causal, witness, irredutibilidade,
intervenções, pluralidade de realizações, constituição P1--P12 e programa inicial. \\

0.2.0 & 09/09/2026 &
Revisão ontológica: define TK-O; introduz perfil constitutivo, ordem explícita de
redução, conjunto de minimais, assinatura intervencional, equivalência causal,
hipótese de Kernel como classe de equivalência, escada de claims e literatura de
enquadramento. \\
\bottomrule
\end{longtable}

\section{Referências}

\begin{thebibliography}{99}

\bibitem{pearl2009}
Pearl, J. (2009).
\emph{Causality: Models, Reasoning, and Inference}.
2nd ed. Cambridge University Press.

\bibitem{woodward2003}
Woodward, J. (2003).
\emph{Making Things Happen: A Theory of Causal Explanation}.
Oxford University Press.
\href{https://doi.org/10.1093/0195155270.001.0001}{doi:10.1093/0195155270.001.0001}.

\bibitem{machamer2000}
Machamer, P.; Darden, L.; Craver, C. F. (2000).
Thinking about mechanisms.
\emph{Philosophy of Science}, 67(1), 1--25.
\href{https://doi.org/10.1086/392759}{doi:10.1086/392759}.

\bibitem{spirtes2001}
Spirtes, P.; Glymour, C.; Scheines, R. (2001).
\emph{Causation, Prediction, and Search}.
2nd ed. MIT Press.
\href{https://doi.org/10.7551/mitpress/1754.001.0001}{doi:10.7551/mitpress/1754.001.0001}.

\bibitem{peters2017}
Peters, J.; Janzing, D.; Schölkopf, B. (2017).
\emph{Elements of Causal Inference: Foundations and Learning Algorithms}.
MIT Press.

\bibitem{rubenstein2017}
Rubenstein, P. K.; Weichwald, S.; Bongers, S.; Mooij, J. M.; Janzing, D.;
Grosse-Wentrup, M.; Schölkopf, B. (2017).
Causal consistency of structural equation models.
\emph{Proceedings of the 33rd Conference on Uncertainty in Artificial Intelligence
(UAI 2017)}, paper 11.

\bibitem{beckers2019}
Beckers, S.; Halpern, J. Y. (2019).
Abstracting causal models.
\emph{Proceedings of the AAAI Conference on Artificial Intelligence}, 33(01),
2678--2685.
\href{https://doi.org/10.1609/aaai.v33i01.33012678}{doi:10.1609/aaai.v33i01.33012678}.

\bibitem{geiger2021}
Geiger, A.; Lu, H.; Icard, T.; Potts, C. (2021).
Causal abstractions of neural networks.
\emph{arXiv:2106.02997}.

\bibitem{iso21838bfo}
ISO/IEC (2021).
\emph{ISO/IEC 21838-2:2021 — Information technology — Top-level ontologies (TLO) —
Part 2: Basic Formal Ontology (BFO)}.
International Organization for Standardization.

\bibitem{obofoundry2021}
Vita, R.; Balhoff, J. P.; et al. (2021).
OBO Foundry in 2021: operationalizing open data principles to evaluate ontologies.
\emph{Database}, 2021, baab069.
\href{https://doi.org/10.1093/database/baab069}{doi:10.1093/database/baab069}.

\end{thebibliography}

\end{document}

